\documentclass[aspectratio=169,11pt]{ctexbeamer}

\usetheme{Madrid}
\usecolortheme{default}
\setbeamertemplate{navigation symbols}{}

\usepackage{graphicx}
\usepackage{booktabs}
\usepackage{array}
\usepackage{xcolor}

\title{Mercury Vibecoding\ \n第二周进展汇报（Team A）}
\author{Team A：Feed \& 同步组}
\institute{课程小组项目}
\date{2026-05-28}

\newcommand{\imgplaceholder}[2]{%
  \begin{center}
    \fcolorbox{gray!60}{gray!8}{%
      \begin{minipage}[c][#2][c]{#1}
        \centering
        \vspace{0.5em}
        {\large 图片占位符}\\[0.4em]
        {\small （后续替换为实际截图）}
      \end{minipage}%
    }
  \end{center}
}

\begin{document}

\begin{frame}
  \titlepage
\end{frame}

\begin{frame}{汇报提纲}
\tableofcontents
\end{frame}

\section{本周目标与完成情况}
\begin{frame}{本周目标（对齐 Team A 职责）}
\begin{itemize}
  \item 完成 Feed 添加/删除与 RSS/Atom 解析
  \item 完成 OPML 导入/导出
  \item 完成 Feed 同步（单源、全量、定时）
  \item 打通文章列表展示闭环
  \item 从本地网页联调升级为 Electron 桌面应用
\end{itemize}
\end{frame}

\begin{frame}{本周交付结果（已完成）}
\begin{itemize}
  \item 数据层：落地 `feeds` / `entries` 表，支持去重写入（`feed_id + guid`）
  \item 解析层：接入 `rss-parser`（RSS/Atom）和 `fast-xml-parser`（OPML）
  \item 服务层：`FeedService`、`OPMLService`、同步流程与错误处理
  \item 前端层：三栏 UI（订阅源 / 文章列表 / 详情）+ 搜索 + 导入导出交互
  \item 桌面化：`Electron Main + Preload + IPC`，前端 API 改为 IPC 优先
\end{itemize}
\end{frame}

\section{实现细节}
\begin{frame}{技术架构（当前可演示版本）}
\begin{center}
\small
用户操作
\ $\downarrow$ \\
Renderer（Vue 3 + TypeScript）
\ $\downarrow$ IPC \\
Main Process（Electron + Node.js）
\ $\downarrow$ \\
Feed/OPML Service
\ $\downarrow$ \\
SQLite（本地数据）
\end{center}

\vspace{0.5em}
\begin{itemize}
  \item 本地优先：无登录、无云端依赖
  \item 跨平台基础：桌面壳已切换为 Electron
\end{itemize}
\end{frame}

\begin{frame}{关键功能演示点}
\begin{columns}[T,onlytextwidth]
\column{0.52\textwidth}
\begin{itemize}
  \item 输入 RSS URL，添加 Feed
  \item 展示文章列表并可搜索
  \item 单源同步/全量同步
  \item OPML 导入（文件/粘贴）
  \item OPML 导出（桌面保存）
\end{itemize}

\column{0.46\textwidth}
\imgplaceholder{0.95\linewidth}{0.48\textheight}
\end{columns}
\end{frame}

\begin{frame}{界面与交互（桌面版）}
\imgplaceholder{0.9\linewidth}{0.58\textheight}
\vspace{0.4em}
\begin{itemize}
  \item 建议替换为：主界面三栏截图（含 Feed、Entry、Detail）
\end{itemize}
\end{frame}

\section{质量与风险}
\begin{frame}{验证结果与当前风险}
\textbf{已验证：}
\begin{itemize}
  \item 后端构建通过，单元测试通过（Feed/Sync/OPML）
  \item 前端构建通过，桌面应用可启动
\end{itemize}

\textbf{风险与待优化：}
\begin{itemize}
  \item 部分中文站点 Feed 源稳定性一般（限流/SSL 兼容）
  \item 前端打包体积偏大（后续可做分包）
  \item 需补充跨平台实机验证（Windows/Linux）
\end{itemize}
\end{frame}

\section{下周计划}
\begin{frame}{下周计划（Team A）}
\begin{itemize}
  \item 完善订阅源异常处理与重试策略
  \item 增加更多 Feed 测试样例与导入容错用例
  \item 协助 Team D 对接打包流程（electron-builder）
  \item 与 Team B/C 联调接口，确保后续功能接入顺畅
\end{itemize}
\end{frame}

\begin{frame}{Q \& A}
\begin{center}
\Large 谢谢大家
\end{center}
\end{frame}

\end{document}
